\documentclass[journal=iecred,manuscript=article]{achemso}
\usepackage{color}
\usepackage{epsf} 
\usepackage{booktabs}
\author{Yang Xiao and Arvind Varma}

\affiliation[Purdue University]
{School of Chemical Engineering, Purdue University, West Lafayette, IN 47907, United States}
\email{avarma@purdue.edu}
\phone{765-494-4075}
\fax{765-494-0805}
\title[An \textsf{achemso} demo]
  {Kinetics of Guaiacol Deoxygenation Using Methane over Pt-Bi Catalyst}

\abbreviations{IR,NMR,UV}
\keywords{Differential reactor, \LaTeX}
\begin{document}
\newpage
TOC graphic 

\begin{center}
\epsfxsize=350pt
\epsffile{TOC.eps}
\end{center}

\newpage
\setlength{\parindent}{0pt}
\begin{center}

\section{Captions for Figures}
\end{center}
Figure 1. Proposed reaction pathways of guaiacol deoxygenation for a) Pt/AC and H$_{2}$ case and b) Pt-Bi/AC and CH$_{4}$ case

Figure 2. Detailed reactions between guaiacol/intermediates and methane

Figure 3. Mass transfer limitation tests for a) internal diffusion effect and b) external diffusion effect

Figure 4. Reaction rate fits for steps R1, R2 and R5 at 400 $^{\circ}$C

Figure 5. A comparison of experimental and predicted flowrate values for a) different catalyst amounts, and b) parity plot

\newpage
\setlength{\parindent}{0pt}
\section{Synopsis}
Using methane as a novel reductant, experimental and modeling efforts related to catalytic deoxygenation kinetics of sustainable bio-oils upgrading are reported.
\section{Abstract}
Using H$_2$ as reductant,
catalytic hydrodeoxygenation (HDO) is typically used for upgrading bio-oils, generally produced from thermal degradation of lignin.
In our recent prior work, methane was used to deoxygenate guaiacol over Pt-Bi catalyst
and was found to exhibit as good deoxygenation performance as hydrogen.
In the present work, using methane as reductant,
detailed reaction pathways and kinetics of guaiacol deoxygenation are studied
using differential and integral operating conditions.
Kinetic parameters including rate constants and activation energies are determined for each individual reaction step.
The model predicted values match experimental data well.
Results from the present work are discussed and compared with the literature values.
The present work provides a practical novel approach for bio-oil upgrading
using methane as reductant instead of hydrogen.

\newpage
\section{Introduction}
Energy security and environmental concerns are drivers for utilization of biomass to produce bio-fuels.
The vast quantity of biomass available in the United States has the potential
 to replace significant amounts of fuels that are currently derived from petroleum sources.
Bio-oil, generally derived from fast pyrolysis of biomass,
is one such candidate fuel. \cite{Tuck2012}
Its production process, however, remains under investigation,
owing primarily to its high oxygen content, 
which may cause instability, poor combustion performance and low heating value. \cite{Nigam2011,Elliott2009,Mortensen2011}
Too high levels of oxygen make pyrolysis oils distinctly different from petroleum fuels.
Oxygen removal from pyrolysis bio-oils, thus, is typically followed as an upgrading process to utilize
these bio-oils as transportation fuels.\cite{Xiu2012, Elliott2015}

Since the pyrolysis bio-oil is always a mixture of more than 400 chemical species,
and different biomass sources result in different bio-oil compositions, model compounds such as phenol, guaiacol (2-methoxyphenol) and benzylphenyl ether\cite{Guvenatam2014},
are typically selected to represent bio-oils for more fundamental studies\cite{Bridgwater2012}
Among these model compounds, guaiacol is attractive because it contains two common oxygenated groups in bio-oils: hydroxyl and methoxyl.
Catalytic hydrodeoxygenation (HDO), as a common bio-oil upgrading approach,
refers to oxygen removal by the use of H$_{2}$ molecules over heterogeneous catalysts.
Using Pt, Pd, Co, etc.\ as primary metals, 
\cite{Wildschut2009} Mo, S, etc.\ as promoters\cite{Wangt2011}, 
and activated carbon (AC), Al$_{2}$O$_{3}$, ZrO$_{2}$, SiO$_{2}$, etc.\ as catalytst supports,\cite{Gutierrez2009,Saidi2014,Wu2013}
various experimental and theoretical studies have been reported for guaiacol catalytic HDO.\cite{Zhao2011, Bui2011, Zhang2013, Mochizuki2014, Olcese2012, Badawi2011, Lee2015}

Although H$_{2}$ is an ideal reductant for oxygen removal from pyrolysis oils,
its use for HDO, however, leads to possible economic concerns owing to high cost for H$_{2}$ production and transportation.
For this reason, we have recently reported a novel bio-oil upgrading approach\cite{Xiao2015},
using CH$_{4}$ instead of  H$_{2}$ as reductant and Pt-Bi/AC as catalyst, to deoxygenate guaiacol efficiently.\
It was found that using the mono-metallic Pt/AC catalyst,
CH$_{4}$ exhibited as good deoxygenation performance as H$_{2}$ in terms of guaiacol conversion and product distribution,
but the catalyst deactivated rapidly.
With addition of bismuth as promoter,
the lifetime of the bi-metallic Pt-Bi/AC was extended significantly without loss of activity,
as compared to the Pt/AC catalyst.
In our other prior works\cite{Gao2014, Gao2015}, using Pt/AC and H$_{2}$, reaction pathways and detailed kinetics of guaiacol deoxygenation process were reported.

On the basis of these prior studies\cite{Xiao2015, Gao2014, Gao2015},
in the present work,
both experimental and modeling investigations were conducted,
to reveal the reaction pathways and kinetics of guaiacol deoxygenation using methane as the reductant,
Kinetics measurements in both differential and integral reactors were carried out, 
and rate constants and activation energy values were determined for the individual reaction steps.
These results are discussed and compared with our prior studies as well as other related works.

\section{Experimental}
\subsection{Materials}
The 5\%Pt-1\%Bi/C catalysts were prepared following the procedure described in our prior works.\cite{Hu2010, Hu2011, Xiao2015}
The metallic precursors were chloroplatinic acid hydrate (99.9\% metal basis) and bismuth chloride (99.999\%), both from Sigma Aldrich.
The 20-120 mesh AC supports were from Norit Americas Inc.
Briefly, Pt and Bi were loaded sequentially using the wet impregnation method. 
First, the Pt and Bi precursors were dissolved in 1.2 mol/L HCl and then added dropwise to the well-stirred AC slurry, 
with stirring continued for at least 8 hrs at room temperature. The mixture was then rinsed and dried in air at 100 $^{\circ}$C before use.
Guaiacol (>98.0\%) and all other calibration compounds (catechol,  phenol and cyclopentanone) were purchased from Sigma-Aldrich. 
GUA, CAT, PHE and CYC are used for abbreviations of Guaiacol, catechol, phenol and cyclopentanone, respectively in the later sections.
Ultra high purity grade gases (99.98\% oxygen, 99.999\% argon, 99.98\% helium and 99.999\% hydrogen) were purchased from Indiana Oxygen.

The 0.5\%Pt/Al$_{2}$O$_{3}$ (metal dispersion=31 $\pm$ 0.5\%) standard (for chemisorption calibration) was from Micromeritics.

\subsection{Kinetics Measurements}
The kinetics measurements were conducted in a fixed-bed reactor setup, described in our prior work. \cite{Gao2014, Gao2015}
Prior to reaction, the packed catalyst was activated at 400 $^{\circ}$C and 1 atm for 4 hours under a gas mixture flow (H$_{2}$:N$_{2}$ = 1:2). 
The reactor was then purged using N$_{2}$ for 30 min.
The standard reactor operating conditions were: 
400 $^{\circ}$C, 1 atm, 0.02-0.05 g catalyst for differential reactions (conversions typically < 6-8\%)
and 0.3-0.5 g catalyst for integral reactions (conversions typically > 20\%), 
total gas (CH$_{4}$:N$_{2}$ = 1:1) flow rate 100 mL/min, 
and guaiacol feed rate 0.025 mL/min (liquid, at room temperature).

The feed flow rates correspond to a molar ratio of 10 between CH$_{4}$ and guaiacol.
Blank tests of AC support with no metal loading were conducted under the standard reaction conditions, and guaiacol conversion was less than 1\%.
Without guaiacol, methane exhibited about 5\% conversion, generating H$_{2}$ and trace C$_{2}$H$_{6}$ on Pt-Bi catalyst under the standard conditions. 
All experiments have carbon mass balances of 92 $\pm$ 2\%, similar to the literature. \cite{Chang2013, Jongerius2013} 
Possible factors affecting mass balance include liquid hold-up in various locations in the system, particularly the condenser, and coke deposit on the catalyst.

\subsection{Product Analysis}
As in our prior works, GC (Agilent GC6890) with flame ionization detector (FID),
equipped with a DB-1701 column (30 m $\times$ 0.25 mm) was used for quantitative analysis of the liquid products. \cite{Gao2014, Gao2015, Xiao2015} 
The gaseous effluent was analyzed using a Micro GC (Agilent 3000A Micro GC) 
equipped with two columns (Column A, MolSieve 5 A, 10 m $\times$ 0.32 mm; 
Column B: Plot U, 8 m $\times$ 0.32 mm) and two thermal conductivity detectors (TCD). 
For the reaction experiments, good repeatability generally within less than 2\% deviation was achieved for all quantitative analysis.

\section{Results and Discussion}
\subsection{Proposed Reaction Pathways}
In our prior work,
when Pt/AC as catalyst and H$_{2}$ as reductant were used for guaiacol deoxygenation,
reactions pathways, as shown in Figure 1a,
were proposed and confirmed by experimental observations\cite{Gao2014}, kinetic modeling and DFT calculations.\cite{Gao2015}
Similar reactions pathways, as shown in Figure 1b, 
apply in the present work since the same product species, including GUA, CAT, PHE and CYC, and similar distribution are found when
Pt-Bi/AC as catalyst and  CH$_{4}$ as reductant are used.\
Both Figure 1a and Figure 1b describe reaction pathways containing five individual steps, denoted as R1-R5,
producing catechol, phenol and cyclopentanone as liquid products,
while different gaseous by-products arise from H$_{2}$ and CH$_{4}$, respectively.
Experimentally, no significant conversion was found when phenol and cyclopentanone were used as reactants in feed, indicating no additional reactions beyond steps R1-R5. 

\begin{center}
\epsfxsize=350pt
\epsffile{Figure01.eps}

Figure 1. Proposed reaction pathways of guaiacol deoxygenation for a) Pt/AC and H$_{2}$ case and b) Pt-Bi/AC and CH$_{4}$ case

\end{center}

Based on the above features, detailed reactions between guaiacol/intermediates and methane are proposed in Figure 2.

\begin{center}
\epsfxsize=350pt
\epsffile{Figure02.eps}

Figure 2. Detailed reactions between guaiacol/intermediates and methane

\end{center}

For a deoxygenation reaction (such as R2 or R3), similar to H$_{2}$,
one molecule of CH$_{4}$ provides only one H atom to oxygen-containing groups,
forming H$_{2}$O or other hydrogenated species,
while the residual CH$_{3}$ combines with another CH$_{3}$ to form a C$_{2}$H$_{6}$ molecule. 

\subsection{Absence of Mass Transfer Limitations}
Before conducting kinetics measurements,
mass transfer limitations were tested in a differential reactor,
using the well-known procedure. \cite{Perego1999}. The experimental results are shown in Figure 3.

\begin{center}
\epsfxsize=500pt
\epsffile{Figure03.eps}

Figure 3. Mass transfer limitation tests for a) internal diffusion effect and b) external diffusion effect
\end{center}

Figure 3a was obtained by varying packed catalyst particle diameter ($d_{p}$)
in the range 100 micron < $d_{p}$ < 1000 micron,
while keeping other operating conditions unchanged.
Figure 3b was obtained by varying feed flow rates,
while maintaining contact time (W/F) and other operating conditions constant.
Figure 3a shows that when catalyst particle diameter ($d_{p}$) is larger than 250 micron,
corresponding to 60 mesh, internal diffusion effects are discerned,
while when ($d_{p}$) < 150 micron, corresponding to 100 mesh, no significant internal effects are found.
As shown in Figure 3b, when feed flow rate is smaller than 1.2 mL/min, 
there is an obvious external mass transfer effect,
while no such effect is observed
when feed flow rate is larger than 1.5 mL/min.
Thus, ($d_{p}$) < 150 micron and feed flow rate larger than 1.5 mL/min were used in the kinetics measurement experiments.
Further, the criteria described by Weisz and Prater \cite{Weisz1954} and
Mears \cite{Mears1971} were both satisfied to verify the absence of mass and heat transfer effects in all experiments.

\subsection{Kinetic Model Development}
The reaction rates of step Ri in Figure 1b
and the corresponding turnover frequencies (TOF) are defined by Eq. (\ref{Eq1}) and Eq. (\ref{Eq2}), respectively
\begin{equation} \label{Eq1}
r_i=\frac{dF}{dW}
\end{equation}
\begin{equation} \label{Eq2}
TOF_i=\frac{dF}{d\left(active~site\right)}=r_i \times \frac{M_{Pt}}{wt\% \cdot D\%}
\end{equation}

Since only one catalyst (Pt-Bi/AC), with specific Pt metal loading amount (wt\%=5\%) and specific Pt metal dispersion (D\%=29\%), was used in the kinetic experiments,
the reaction rate and TOF are inter-converted
as indicated in Eq. (\ref{Eq1}) and Eq. (\ref{Eq2}).

The reaction pathways, including five individual steps (R1-R5) as shown in Figure 1b,
were established under integral operating conditions.
When using guaiacol as feed, under differential operating conditions
(guaiacol conversion less than 6-8\%),
catechol is the only product generated, i.e.\ only step R1 prevails
while steps R3 and R4 are negligible. 
When catechol is used as feed,
also under differential conditions, both phenol and cyclopentanone (corresponding to steps R2 and R5) exist in the liquid products.
These experimental observations imply that when feeding guaiacol under differential operating conditions,
R1 is faster than R3 and R4, and phenol and cyclopentanone are not produced owing to low concentration of catechol.
Therefore, kinetic parameters for steps R1, R2 and R5 could be obtained by running differential reactions.
Using kinetic parameters of R1, R2 and R5,
those for R3 and R4 could be further regressed by running integral reactions.

When differential operating conditions apply, reaction rates are obtained by Eq. (\ref{Eq3}).
\begin{equation} \label{Eq3}
r_i=F_{0} \times \frac{X}{W}
\end{equation}

Based on the design equation for a fixed-bed reactor and the proposed pathways in Figure 1b, 
consumption/formation rates of guaiacol, catechol, phenol and cyclopentanone are shown in Eqs. (\ref{Eq4} -~\ref{Eq7}).

\begin{equation} \label{Eq4}
\frac{dF_{GUA}}{d_{W}}=-r_1-r_3-r_4
\end{equation}
\begin{equation} \label{Eq5}
\frac{dF_{CAT}}{d_{W}}=r_1-r_2-r_5
\end{equation}
\begin{equation} \label{Eq6}
\frac{dF_{CYC}}{d_{W}}=r_4+r_5
\end{equation}
\begin{equation} \label{Eq7}
\frac{dF_{PHE}}{d_{W}}=r_2+r_3
\end{equation}

An overall material balance yields the following Eqs. (~\ref{Eq8} -~\ref{Eq10}).

\begin{equation} \label{Eq8}
F_{GUA}+F_{CAT}+F_{CYC}+F_{PHE}=F_{0}
\end{equation}
%\begin{equation} \label{Eq9}
%P_{GUA}+P_{CAT}+P_{CYC}+P_{PHE}=P_{0}
%\end{equation}
\begin{equation} \label{Eq10}
P_{i}=F_{i} \times \frac{P_{0}}{F_{0}}
\end{equation}

As described in the Kinetics Measurements section,
excess methane (molar ratio of methane to reactants is $\sim$10:1)
is used in all cases,
thus partial pressure of methane may be considered as a constant contribution to the reaction rate expressions.
The following data fitting and optimization are based on experimental measurements at various temperatures from 325-450 $^{\circ}$C.
Although only results at 400 $^{\circ}$ C are shown below because this standard operating temperature
gives optimum deoxygenation performance of guaiacol \cite{Xiao2015}
and results at 400 $^{\circ}$ C were repeated more than 5 times,
fitting and optimization results for different temperature give similar conclusions.
All experiments at other temperatures were repeated at least twice and are discussed in the section Activation Energies and Temperature Effect.

\subsection{Differential Reactions for Steps R1, R2 and R5}
When differential operating conditions apply at 400 $^{\circ}$C, for steps R1, R2 and R5,
the corresponding reaction rates $r_1$, $r_2$ and $r_5$ were obtained by Eq. (\ref{Eq3}).

\begin{center}
\epsfxsize=400pt
\epsffile{Figure04.eps}

Figure 4. Reaction rate fits for steps R1, R2 and R5 at 400 $^{\circ}$C
\end{center}

Using different feed partial pressures of guaiacol or catechol, 
reaction rates were measured and plotted
with respect to average bed reactant partial pressures, as shown in Figure 4.
All three data sets follow linear regressions,
indicating first orders for steps R1, R2 and R5.
Thus, the corresponding reaction rates at 400 $^{\circ}$C are described by in Eqs. (\ref{Eq11} -~\ref{Eq13}),
where the regressed values of $k_1$, $k_2$ and $k_5$ are listed in Table 1.
The reaction rate regressions of R1, R2 and R5 at other temperatures are also essentially linear,
as discussed later in the section Activation Energies and Temperature Effect.

\begin{equation} \label{Eq11}
r_1=k_1 P_{GUA}
\end{equation}
\begin{equation} \label{Eq12}
r_2=k_2 P_{CAT}
\end{equation}
\begin{equation} \label{Eq13}
r_5=k_5 P_{CAT}
\end{equation}

\begin{center}

Table 1. Reaction rate constants for steps R1, R2 and R5 at 400 $^{\circ}$C;
units of $k_i$ $mol/(g_{cat} \cdot s \cdot atm)$

\begin{tabular}{c c c c c c c c c c c c c}
\toprule
           Steps &  $k_i\times 10^{-5}$\\
\midrule
          R1 & 8.29\\
          R2 & 9.71\\
          R5 & 4.47 \\
\bottomrule
\end{tabular}
\end{center}

Using H$_2$ as reductant, Runnebaum et al.\ \cite{Runnebaum2012} also reported 1st order kinetics for step R1 over Pt/Al${_2}  $O${_3}$ catalyst, 
while our prior work \cite{Gao2015} fitted a 2nd order for the same step.
The difference between the present work and our prior work may be due to different reductants (H$_2$ or CH$_4$).
For steps R2 and R5, both the present work and our prior work demonstrate 1st order kinetics.
The reaction rate constant values, however, are smaller when CH$_4$ is used as reductant, indicating a higher activation energy for CH$_4$.

\subsection{Integral Reactions of Steps R3 and R4}

As shown above, under differential operating conditions,
step R1 is decomposed from the reaction network (Figure 1) when feeding guaiacol over the catalyst.
Similarly, steps R2 and R5 are decomposed under differential operating conditions when feeding catechol.
To obtain insight into steps R3 and R4,
integral operating conditions were used while feeding guaiacol
by packing more catalyst,
producing distinct amounts of phenol and cyclopentanone.
These, obviously, introduce not only steps R3 and R4, but also steps R1, R2 and R5.
Thus Eqs. (\ref{Eq4} -~\ref{Eq7}) can be re-written as Eqs. (\ref{Eq22} -~\ref{Eq25}),
where $k_i$ and $n_i$ are reaction rate constants and orders for steps $R_i$.

\begin{equation} \label{Eq22}
\frac{dF_{GUA}}{d_{W}}=-r_1-r_3-r_4=-k_1 P_{GUA}^{n_1}-k_3 P_{GUA}^{n_3}-k_4 P_{GUA}^{n_4}
\end{equation}
\begin{equation} \label{Eq23}
\frac{dF_{CAT}}{d_{W}}=r_1-r_2-r_5=-k_1 P_{GUA}^{n_1}-k_2 P_{CAT}^{n_2}-k_5 P_{CAT}^{n_5}
\end{equation}
\begin{equation} \label{Eq24}
\frac{dF_{CYC}}{d_{W}}=r_4+r_5=k_4 P_{GUA}^{n_4}+k_5 P_{CAT}^{n_5}
\end{equation}
\begin{equation} \label{Eq25}
\frac{dF_{PHE}}{d_{W}}=r_2+r_3=k_2 P_{CAT}^{n_2}+k_3 P_{GUA}^{n_3}
\end{equation}

In this description, the reaction orders for steps 1, 2 and 5 are also kept flexible. With $k_i$ and $n_i$ as variables (total 10),
predicted values of flowrates for GUA, CAT, CYC and PHE can be calculated
by Runge-Kutta numerical integration using the following initial conditions:

\begin{equation} \label{Eq18}
F_{GUA}=F_0,~F_{CAT}=0,~F_{PHE}=0~and~F_{CYC}=0,~when~W=0.
\end{equation}

In general, for parameter estimation involving exponents,
Weighted Least-Squares method is recommended and was followed in the present work\cite{Norris1981, Fogler2005, Box2005}.   
This involves minimizing the error estimate given by Eq.~(\ref{Eq19}). 

\begin{equation} \label{Eq19}
\sigma^2=\sum_{i}{w_i(F_{exp_i}-F_{pred_i})^2},
\end{equation}

where $~i=GUA,~CAT,~PHE~and~CYC$ if integral conditions apply,
while $~i=GUA~and~CAT$ if differential conditions apply,
and $w_i$ are the weighting functions.
In principle, any positive weighting function can be used.
For single reactions, the appropriate $w_i$ have been reported by de Levie \cite{Levie1986},
but are not available for simultaneous multi-reaction network, as in the present case.
Following de Levie, $w_i=1$, $F^2_i$, $F^4_i$ and $F^6_i$ are
appropriate for single reactions with orders 0, 1, 2 and 3, respectively. \cite{Levie1986}

Since the differential reaction analysis yielded first-order for steps R1, R2 and R5, 
we used $w_i=F^2_i$ as the starting weighting functions.
A sequential optimization strategy was then employed where only
2 ($k_3$ and $k_4$) or 4 ($k_3$, $k_4$, $n_3$ and $n_4$) variables out of 10 were kept flexible at the initial stage.
Using these results as initial guesses,
all 10 variables were allowed to vary in the final stage
and the corresponding optimized parameters are shown in Table 2. Since the reaction orders $n_3$ and $n_4$ are 2, 
$w_i=F^4_i$ was also tested, and converged to the same $n_i$
as for $w_i=F^2_i$ and essentially the same $k_i$ values as well.
Thus it appears that $w_i=F^2_i$ is an appropriate choice for the weighting function.

\begin{center}
Table 2. Optimized $n_i$ and $k_i$ values for $w_i=F^2_i$;
units of $k_i$ $mol/(g_{cat} \cdot s \cdot atm^{n_i})$ 
\begin{tabular}{c c c c c c c c c c c c c}
\toprule
 $n_1$ & $n_2$ & $n_3$ & $n_4$ & $n_5$ & $k_1\times 10^{-5}$ & $k_2\times 10^{-5}$ & $k_3\times 10^{-5}$ & $k_4\times 10^{-5}$ & $k_5\times 10^{-5}$\\
\midrule
 1.00 & 1.00 & 2.00 & 2.00 & 1.00 & 8.38 & 9.68 & 23.8  & 1.76 & 4.46\\
\bottomrule
\end{tabular}
\end{center}

Using data from both integral and differential operating conditions,
a comparison of the experimental and predicted flowrates is shown in Figure 5.
Figure 5a shows that predicted flowrates are close to the experimental values for various catalyst amounts.
Figure 5b summarizes the goodness-of-fit in a parity plot.
The values for all species are close to the diagonal line and relatively evenly distributed on both sides, indicating a good fit (normalized RMS error = 2.75\%).

\begin{center}
\epsfxsize=500pt
\epsffile{Figure05.eps}

Figure 5. A comparison of experimental and predicted flowrate values for a) different catalyst amounts, and b) parity plot
\end{center}

\subsection{Activation Energies and Temperature Effect}
In the above sections, all reaction orders and rate constants at 400 $^{\circ}$C were determined.
Using the Arrhenius equation (\ref{Eq31}), reaction activation energies ($E_a$) and exponential factors ($A$) were calculated.

\begin{equation} \label{Eq31}
k=Aexp\left(-\frac{E_a}{RT}\right)
\end{equation}

The $E_a$ values are listed in Table 3 and also compared with our prior work \cite{Gao2015},
where hydrogen and Pt/AC were used.

\begin{center}
Table 3. Reaction activation energy values, along with comparison for CH$_4$ and H$_2$ cases

\begin{tabular}{c c c c c c c c c c c c c}
\toprule
$E_a,~kJ/mol$ & $R_1$ & $R_2$ & $R_3$ & $R_4$ & $R_5$\\
\midrule
$CH_4~case$ & 146 & 110 & 166 & 184 & 141\\
$H_2~case$ & 126 & 100 & 93 & 149 & 125\\
\bottomrule
\end{tabular}
\end{center}

The activation energy of step R1 was reported to be 71.2, 58.7 and 89.1 $kJ/mol$
for Co$-$Mo, Ni$-$Mo and Ni$-$Cu catalysts, respectively, \cite{Laurent1994,Bykova2013}
when H$_2$ was used as reductant.
There is no data available for the CH$_4$ case in the literature.
Using Pt-based catalysts, 
our activation energies of step R1 are 146 and 126 $kJ/mol$ for the CH$_4$ and H$_2$ cases, respectively, both higher than literature data.
The difference likely owes to the catalyst itself,
including Pt metal nature, Pt particle size, support effect and Bi promoter addition \cite{Deepa2014, Xiao2015}.
Table 3 also shows that activation energies for the CH$_4$ case are always 10-35 $kJ/mol$ higher (except R3) than those for the H$_2$ case. 
This is an explanation for the fact that, to achieve a similar guaiacol conversion at 300 $^{\circ}$C for the H$_2$ case,
400 $^{\circ}$C is needed for the CH$_4$ case as shown in our prior work.\cite{Xiao2015}.
For step R3, activation energy of the H$_2$ case is only 93 $kJ/mol$,
while it is 166 $kJ/mol$ for the CH$_4$ case.
This is likely because H$_2$ promotes step R3 by forming a more stable by-product CH$_3$OH than C$_2$H$_6$, the by-product from the CH$_4$ case.

In the present work, kinetics of guaiacol deoxygenation using methane as reductant over Pt-Bi/AC catalyst is 
studied by both experimental and modeling approaches.
It is shown that the liquid phase products remain the same as for the case using hydrogen as reductant.\cite{Gao2015}
The activation energies and other kinetic parameters are obtained based on data fit.
The experimental data matches model predictions well.
Modeling results are discussed, analyzed and compared with literature reports.
The present work provides a new practical approach for bio-oil upgrading
using methane as reductant instead of hydrogen.

\section{Nomenclature}

\begin{tabbing}

01\hspace*{50bp} \=02\hspace*{90bp} \=03\hspace*{50bp}\kill

$A_i$  \> $[mol\left(s \cdot g_{cat}\right)]$   \> Pre-exponential factor of step i \\
$D$  \> $[\%]$   \> Metal dispersion\\
$Ea_i$  \> $[kJ/mol]$   \> Activation energy of step i \\
$F_0$  \> $[mol/h]$   \> Feed flowrate \\
$F_i$  \> $[mol/h]$   \> Flowrate of species i \\
$k_i$  \> $[mol/(g_{cat} \cdot s \cdot atm^n) ]$   \> Reaction rate constant of step i when the reaction order is n\\
$M_{Pt}$  \> $[kg/kmol]$   \> Molecular weight of Pt \\
$N_i$  \> $[mol]$   \> Amount of species i \\
$n_i$  \> $[mol]$   \> Reaction order of step  i \\
$P_i$  \> $[atm]$   \> Partial pressure of species i \\
$R$  \> $[J/\left(K \cdot mol \right)]$   \> Gas constant \\
$r_i$  \> $[mol\left(s \cdot g_{cat}\right)]$  \> Reaction rate of step i \\
$t$  \> $[s]$   \> Contact time \\
$TOF_i$  \> $[s^{-1}]$  \> Turnover frequency of step i \\
$W$  \> $[g]$   \> Catalyst packing amount\\
$W_i$  \> $[-]$   \> weighting function for species i\\
$wt$  \> $[\%]$   \> Metal loading percentage \\
\end{tabbing}

\section{Acknowledgments}
This work was supported by the R. Games Slayter Fund.
The authors thank Ms. Marina Lourencone for her valuable assistance.

\bibliography{Kinetics_CH4}
\end{document}
